%% kapa: exam=1 round=27 year=2016 period=2 form=A subject=accounting subject_ko=회계학 seq=36 number=76
%% source: Q-Net 감정평가사 기출문제 (cst00309) artlSeq=5169402 file=extracted/1008/5169402/제27회 감정평가사 1차 문제/제27회 감정평가사 2교시문제(A형).pdf
%% conversion: kordoc 4.24.51 -> tools/gen1.py (auto); review=proofread-sonnet; flags=none
\begin{question}{76}
(주)감평은 A, B 두 개의 사업부만 두고 있다. 투자수익률과 잔여이익을 이용하여 사업부를 평가할 때 관련 설명으로 옳은 것은? (단, 최저필수수익률은 6\%라고 가정한다.)

{\small\centering
\begin{tabular}{|c|r|r|}\hline
구 분 & \multicolumn{1}{c|}{A사업부} & \multicolumn{1}{c|}{B사업부} \\ \hline
투자금액 & ￦250,000,000 & ￦300,000,000 \\ \hline
감가상각비 & 25,000,000 & 28,000,000 \\ \hline
영업이익 & 20,000,000 & 22,500,000 \\ \hline
\end{tabular}\par}
\begin{choices}
  \item A사업부와 B사업부의 성과는 동일하다.
  \item A사업부가 투자수익률로 평가하든 잔여이익으로 평가하든 더 우수하다.
  \item B사업부가 투자수익률로 평가하든 잔여이익으로 평가하든 더 우수하다.
  \item 투자수익률로 평가하는 경우 B사업부, 잔여이익으로 평가하는 경우 A사업부가 각각 더 우수하다.
  \item 투자수익률로 평가하는 경우 A사업부, 잔여이익으로 평가하는 경우 B사업부가 각각 더 우수하다.
\end{choices}
\end{question}
\answer{2}
